
\paragraph{Evidence and limits.}
This is a paper proof, accompanied by exact rational finite identity checks.
It is not Lean formalized or externally peer reviewed. No novelty claim is
made: the argument uses polynomial moment matching. It rules out a bounded
local-depth pilot strategy, including arbitrary feedback between coins. It
does not resolve the unrestricted rare-event minimax problem.

\paragraph{Experiment.}
Fix $0<a<1/2$ and the known law $\mu(dd)=2d\,dd/a^2$ on $(0,a]$.
The sign allocation $q$ is unknown, $0\le q\le1$ and $\int q\,d\mu\le\rho$,
where $0<\rho\le1/4$ is known. Each iid fresh coin has hidden $D\sim\mu$,
hidden $S\mid D=d\sim\mathrm{Ber}(q(d))$, and supplies conditionally iid
$\mathrm{Ber}(1/2+(2S-1)D)$ observations. An algorithm may use arbitrary
independent randomization, create, revisit, interleave and abandon coins, and
let all these decisions depend on all earlier observations. Only flips and
identities are observed. Every flip is charged. The restriction considered
below is solely that each individual coin supplies at most $2K$ flips.

\begin{lemma}[An elementary bounded polynomial]\label{lem:poly}
For integer $K\ge1$, let $P_n$ be the degree-$n$ Legendre polynomial defined
by Rodrigues' formula, $n=2K+1$, and put
\[
 c=P_n'(0),\qquad u_K(x)=P_n(x)/(cx),\quad u_K(0)=1.
\]
Then $u_K$ is an even polynomial of degree $2K$ satisfying
\[
 |u_K(x)|\le1\quad(0\le x\le1),\qquad
 \int_0^1x^{2j+2}u_K(x)\,dx=0\quad(0\le j<K),
\]
and
\[
 \int_0^1 2x u_K(x)\,dx=\frac1{(K+1)(2K+1)}=:m_K.
\]
\end{lemma}
\begin{proof}
For completeness the polynomial facts needed here follow from
\[
 P_n(x)=\frac1{2^n n!}\frac{d^n}{dx^n}(x^2-1)^n.
\]
Expansion gives
\[
 c=(-1)^K(2K+1){\binom{2K}{K}}4^{-K},\qquad
 P_{2K}(0)=(-1)^K\binom{2K}{K}4^{-K}=:v_K.
\]
It also gives the differential and derivative identities
\[
 (1-x^2)P_n''-2xP_n'+n(n+1)P_n=0,
 \qquad P_{n+1}'-P_{n-1}'=(2n+1)P_n.
\]
These identities can equivalently be verified coefficient by coefficient
from the displayed expansion. Rodrigues' formula at $x=1$ gives $P_j(1)=1$:
in the derivative of $(x-1)^j(x+1)^j$ only the term differentiating the first
factor $j$ times survives at $1$.

Substitute $P_n=cxu_K$ in the differential equation and set
$\lambda=n(n+1)-2=2K(2K+3)>0$. On $(0,1)$,
\[
 (1-x^2)u_K''+(2/x-4x)u_K'+\lambda u_K=0.
\]
The energy $H(x)=u_K(x)^2+(1-x^2)u_K'(x)^2/\lambda$ satisfies
\[
 H'(x)=(6x-4/x)u_K'(x)^2/\lambda.
\]
Thus $H$ decreases up to $\sqrt{2/3}$ and increases thereafter, and its
maximum on $[0,1]$ is at an endpoint. Since $u_K$ is even, $H(0)=1$;
$H(1)=1/c^2\le1$. The last inequality follows from $|c|\ge1$: the value
at $K=0$ is $1$ and the ratio of successive absolute values is
$(2K+3)/(2K+2)>1$. Since $H\ge u_K^2$, the asserted bound follows.

Integrating Rodrigues' formula by parts $n$ times shows that $P_n$ is
orthogonal on $[-1,1]$ to every polynomial of degree less than $n$; the
boundary terms vanish because $(x^2-1)^n$ has zeros of order $n$ at both
endpoints. Therefore
\[
 \int_0^1x^{2j+2}u_K(x)\,dx
 =\frac1{2c}\int_{-1}^1x^{2j+1}P_n(x)\,dx=0.
\]
Finally $v_{K+1}=-(2K+1)v_K/(2K+2)$ and the derivative identity give
\[
 \int_0^1P_{2K+1}(x)\,dx
 =\frac{v_K-v_{K+1}}{4K+3}=\frac{v_K}{2K+2}.
\]
Multiply by $2/c$, where $c=(2K+1)v_K$, to obtain $m_K$.
\end{proof}

\begin{theorem}[Feedback does not repair bounded-depth nonidentifiability]
\label{thm:cap}
For any $K\ge1$, define $h=a\sqrt\rho$ and
\[
 q_\pm(d)=\frac{1\pm u_K(d/h)}2\,\mathbf1_{\{d\le h\}}.
\]
Both are admissible allocations, with
\[
 \theta_\pm=\frac\rho2(1\pm m_K),\qquad
 \theta_+-\theta_-=\frac\rho{(K+1)(2K+1)}.
\]
Every population algorithm obeying the per-coin cap $2K$ has identical
finite-prefix transcript distributions under these two allocations,
including arbitrary cross-coin feedback and pilot-dependent acquisition.
If its terminal output is measurable and the procedure terminates almost
surely under both laws, its output distributions also coincide. Consequently
its worst-case RMS error is at least
\[
 \frac\rho{2(K+1)(2K+1)},
\]
regardless of its expected total work or number of coins.
\end{theorem}
\begin{proof}
Lemma~\ref{lem:poly} implies $0\le q_\pm\le1$. A change of variable
$d=hx$ gives the displayed targets, each at most $\rho$. This uses only
the global rarity constraint: the pointwise allocations can be as large as
one in the small difficult population.

For a particular bit word $w$ of length $r$ and head count $s$, its likelihood
conditional on a positive or negative sign is
\[
 \ell_\pm(d)=(1/2\pm d)^s(1/2\mp d)^{r-s}.
\]
The difference $\ell_+(d)-\ell_-(d)$ is an odd polynomial of degree at most
$r$. When $r\le2K$, all its monomials have degrees $2j+1$ with $j<K$.
For each such degree,
\[
 \int(q_+-q_-)(d)d^{2j+1}\,\mu(dd)
 =\frac{2h^{2j+3}}{a^2}\int_0^1u_K(x)x^{2j+2}\,dx=0.
\]
Hence the probability of every local word of length at most $2K$ is exactly
the same under the two allocations.

Condition on the algorithm's independent seed. For any finite global
transcript, its probability is the product of the probabilities of the
local words of the distinct coins it sampled, multiplied by the indicator
that its scheduling and stopping decisions agree with the transcript.
The indicator is parameter independent; every local word has length at most
$2K$. Thus each finite prefix has the same probability under the two laws.
Integrating the seed preserves equality. Coins may be indexed by their
first charged observation, so free unused identities cause no difficulty.
Equivalently, equality follows successively from identical conditional
next-bit probabilities for every eligible local history.

For an almost surely terminating procedure, a terminal event intersected
with the event of termination within $N$ observations belongs to a finite
prefix. Its probabilities coincide for every $N$ and converge to the full
terminal probabilities under both laws. Finally, if $Z$ has this common
output law and $\Delta=\theta_+-\theta_-$, then pointwise
\[
 \tfrac12[(Z-\theta_+)^2+(Z-\theta_-)^2]
 =(Z-(\theta_++\theta_-)/2)^2+\Delta^2/4\ge\Delta^2/4.
\]
Taking expectations, including infinite second moments, proves the RMS bound.
\end{proof}

\begin{theorem}[A depth-tail necessity for unrestricted feedback]\label{thm:tail}
Keep the two allocations and target gap $\Delta_K=\rho m_K$ from
Theorem~\ref{thm:cap}. Now allow an unrestricted population algorithm,
including arbitrarily many observations on a coin, and assume almost-sure
termination under these two laws. Let $E_K$ be the event that some individual
coin receives more than $2K$ flips. Then
\[
 p_K:=\mathbb P_+(E_K)=\mathbb P_-(E_K).
\]
If the RMS error is at most $\varepsilon$ under both laws, then
\[
 p_K\ge 1-\frac{4\varepsilon^2}{\Delta_K^2}.
\]
In particular, $\Delta_K\ge4\varepsilon$ forces $p_K\ge3/4$ under each law.
Moreover, separately under both laws,
\[
 \mathbb E_\pm T\ge(2K+1)p_K.
\]
The right side of the probability lower bound is informative only when
positive; it can always be replaced by its maximum with zero.
\end{theorem}
\begin{proof}
Construct a capped experiment from the given algorithm. Follow all its
decisions until it terminates, or until it is about to request the
$(2K+1)$st observation of some coin. In the latter case stop immediately
\emph{before} that observation and record an exit flag. Otherwise record
the original terminal output and a no-exit flag. This auxiliary experiment
obeys the local cap $2K$. It terminates almost surely under both laws: on
paths with no exit it follows the original terminating procedure, and on
exit paths it stops earlier. Theorem~\ref{thm:cap} therefore gives equality
of its terminal laws. In particular the exit probabilities agree, proving
the first assertion. This equality concerns the event of the first excessive
query, whose decision is made using only earlier allowed observations; it
does not assert equality of later unrestricted transcripts.

Let $\nu$ be the subprobability measure of the original output together with
the no-exit event. The capped experiment shows that it is the same under
both laws and has total mass $1-p_K$. If $Z$ denotes the output coordinate,
nonnegativity of squared errors gives
\begin{align*}
 \mathbb E_+(Z-\theta_+)^2+\mathbb E_-(Z-\theta_-)^2
 &\ge\int\bigl[(z-\theta_+)^2+(z-\theta_-)^2\bigr]\,\nu(dz)\\
 &\ge (1-p_K)\Delta_K^2/2.
\end{align*}
The assumed RMS bounds make the left side at most $2\varepsilon^2$, giving
the probability claim. This common-submeasure argument also gives total
variation at most $p_K$ between the unrestricted output laws; there is no
factor two, since their common part has mass $1-p_K$.
On the exit event the original algorithm actually requests and obtains a
coin's $(2K+1)$st flip, so its total charged work is at least $2K+1$.
Taking expectations proves the work bound, even when the expectations are
infinite.
\end{proof}

\paragraph{The first nontrivial calibration example.}
For $K=2$,
\[
 u_2(x)=1-\frac{14}3x^2+\frac{21}5x^4,\quad m_2=1/15,
\]
and the targets are $8\rho/15$ and $7\rho/15$. No population policy with
at most four flips per coin can have uniform RMS error below $\rho/30$,
even with arbitrarily many pilot and production coins and nonlinear output.
The two profiles agree on every pilot statistic measurable from these
observations. A pilot that fits one of the profiles and uses that fit for
later calibration therefore cannot validate its target uniformly.

\paragraph{What remains open.}
This theorem permits arbitrary feedback but imposes a fixed local cap. The
main expected-budget experiment has no such cap. The depth-tail corollary
nevertheless applies to that unrestricted experiment: sufficiently accurate
procedures must enter greater local depths with a quantified probability.
The target gap decreases quadratically with the cap and proves neither an iterated logarithm in the
unrestricted rate nor the optimality of any proposed estimator. An attempt
to turn pilot-learned survivor weights into a uniform estimator remains
unsuccessful: calibration requires controlling the response-weighted error
of the nuisance fit, and finite-depth agreement does not do that.
